\chapter{Redes de Computadores}
\label{cap:redes-teoria}

Este capítulo cobre os fundamentos de redes que aparecem na prova, com a
mecânica por trás de cada conceito --- não só o rótulo. A versão resumida,
com o recorte do que realmente cai no Perfil 3, está na
\texttt{apostila/}; aqui o objetivo é entender \textbf{por que} a pilha de
protocolos funciona do jeito que funciona.

\section{Modelos de referência: OSI e TCP/IP}

\subsection{Encapsulamento: a ideia que sustenta o modelo em camadas}

\begin{conceito}[Camada é divisão de trabalho, não gaveta de classificação]
O problema de fazer dois programas em máquinas diferentes conversarem é
grande demais para resolver de uma vez --- envolve sinal elétrico, entrega
ao vizinho físico, roteamento entre redes, garantia de entrega, e o
significado dos dados em si. O modelo em camadas \textbf{divide esse
problema} em fatias menores, cada uma resolvida por uma camada diferente,
que não precisa saber como as outras resolvem o próprio pedaço.

O mecanismo que viabiliza isso é o \textbf{encapsulamento}: ao descer a
pilha (enviando), cada camada recebe os dados da camada de cima, faz seu
trabalho e \textbf{acrescenta um cabeçalho próprio} antes de passar adiante.
Ao subir a pilha (recebendo), cada camada \textbf{remove} o cabeçalho que
corresponde a ela e entrega o resto para cima. É como uma carta dentro de
um envelope, dentro de uma caixa, dentro de outra caixa --- cada camada só
abre a própria embalagem.
\end{conceito}

\begin{conceito}[OSI --- 7 camadas]
\begin{center}
\begin{tabular}{@{}p{0.6cm}p{2.2cm}p{4.4cm}p{3.4cm}@{}}
\toprule
\# & \textbf{Camada} & \textbf{Problema que resolve} & \textbf{Exemplos} \\
\midrule
7 & Aplicação & serviço que o usuário/programa efetivamente quer & HTTP,
DNS, SMTP \\
6 & Apresentação & formato dos dados --- tradução, cifra, compressão &
codificação, TLS (discutível) \\
5 & Sessão & controle do diálogo --- abrir, manter, retomar de um
\textit{checkpoint} & --- \\
4 & Transporte & entregar ao \textbf{processo} certo (porta), com ou sem
garantia de ordem/entrega & TCP, UDP \\
3 & Rede & achar o caminho \textbf{entre redes diferentes} & IP, roteador \\
2 & Enlace & entregar ao vizinho \textbf{dentro do mesmo segmento} físico
(endereço MAC) & switch, Ethernet \\
1 & Física & converter bits em sinal no meio físico & cabo, rádio \\
\bottomrule
\end{tabular}
\end{center}

Cada camada tem uma unidade de dados própria (PDU), e o nome dessa unidade
denuncia a camada em qualquer questão: \textbf{segmento} (TCP) na
Transporte, \textbf{pacote} na Rede, \textbf{quadro} na Enlace, \textbf{bit}
na Física.
\end{conceito}

\begin{conceito}[TCP/IP --- 4 camadas, a pilha realmente implementada]
O modelo OSI é o \textbf{referencial de ensino} (didático, nunca foi
implementado à risca como está descrito); a pilha que roda de verdade na
Internet é a \textbf{TCP/IP}, mais enxuta: Acesso à Rede (funde OSI 1--2),
Internet (equivale a OSI 3), Transporte (equivale a OSI 4), Aplicação (funde
OSI 5--7). Quando um enunciado menciona “camada de sessão” ou “camada de
apresentação”, está necessariamente falando do modelo \textbf{OSI} --- essas
duas não têm correspondência separada no TCP/IP.
\end{conceito}

\begin{exemplo}[Identificando a camada pela função descrita, não pelo nome]
“Um protocolo insere marcadores no fluxo de dados para permitir que, se a
conexão cair no meio de uma transferência grande, ela retome do ponto
interrompido em vez de recomeçar do zero.” $\to$ isso é controle de
\textbf{diálogo com checkpoint} $\to$ camada de \textbf{Sessão} (OSI 5) ---
não é Transporte (que garante entrega e ordem, mas não gerencia retomada de
diálogo em alto nível) nem Apresentação (que trata de formato dos dados).

“Um roteador recebe um pacote e decide por qual próximo salto encaminhá-lo,
com base no endereço de destino.” $\to$ camada de \textbf{Rede} (OSI 3) ---
é sobre achar caminho entre redes.

“Um switch recebe um quadro e decide para qual porta física repassá-lo, com
base na tabela de endereços MAC aprendidos.” $\to$ camada de \textbf{Enlace}
(OSI 2) --- a decisão fica dentro do mesmo segmento físico.
\end{exemplo}

\begin{regrapratica}[Roteiro para identificar a camada em qualquer cenário]
Pergunte, nesta ordem: o problema é \textbf{entre duas máquinas do mesmo
segmento físico}? Enlace. É \textbf{entre redes diferentes} (precisa
escolher rota)? Rede. Chegou \textbf{à máquina certa mas falta garantir
processo/porta, ordem ou retransmissão}? Transporte. É sobre
\textbf{retomar/gerenciar um diálogo} (checkpoint, sessão)? Sessão. É sobre
\textbf{formato, cifra ou compressão} dos dados em si? Apresentação. É o
\textbf{serviço final} que o usuário/programa pediu? Aplicação. A banca
sempre descreve o comportamento, nunca pergunta “qual é o número da camada
X” isoladamente --- o número é consequência de identificar a função.
\end{regrapratica}

\section{Equipamentos de rede por camada}

\begin{conceito}[Hub, switch e roteador --- a inteligência cresce com a camada]
\begin{itemize}
\item \textbf{Hub} (camada 1 --- Física): recebe um sinal em uma porta e
\textbf{repete para todas as outras}, sem entender endereço nenhum. É “burro”
por definição --- não sabe distinguir destinatários, então gera colisão e
desperdiça banda.
\item \textbf{Switch} (camada 2 --- Enlace): aprende quais endereços
\textbf{MAC} estão em quais portas (observando o tráfego que passa) e
encaminha cada quadro \textbf{só para a porta certa} --- elimina a
repetição cega do hub.
\item \textbf{Roteador} (camada 3 --- Rede): usa o endereço \textbf{IP} de
destino para decidir por qual interface encaminhar um pacote, inclusive
\textbf{entre redes diferentes} --- é o único dos três capaz de interligar
duas redes distintas.
\end{itemize}
Existe switch de camada 3 (que também roteia), mas o par cobrado com mais
frequência é a associação direta \textbf{switch = camada 2}, \textbf{roteador
= camada 3}.
\end{conceito}

\section{TCP $\times$ UDP}

\begin{conceito}[Os dois protocolos de transporte resolvem trade-offs opostos]
Ambos entregam dados a um \textbf{processo} (via porta), mas fazem
concessões diferentes entre confiabilidade e velocidade:

\begin{center}
\begin{tabular}{@{}p{2.6cm}p{4.6cm}p{4.6cm}@{}}
\toprule
& \textbf{TCP} & \textbf{UDP} \\
\midrule
Conexão & orientado a conexão (estabelece com \textit{handshake} de 3 vias
antes de trocar dados) & sem conexão --- envia direto \\
Confiabilidade & garante entrega e ordem (retransmite o que se perde,
reordena o que chega fora de ordem) & não garante nada disso --- se um
pacote se perde, não há retransmissão automática \\
Custo & mais overhead (cabeçalho maior, controle de estado da conexão) &
mais leve e rápido \\
Uso típico & web, e-mail, transferência de arquivo --- onde perder um byte
é inaceitável & streaming, DNS, VoIP --- onde chegar rápido importa mais
que chegar perfeito (um pacote perdido de voz não vale a pena retransmitir,
já ficou velho) \\
\bottomrule
\end{tabular}
\end{center}
\end{conceito}

\begin{conceito}[O \textit{three-way handshake} do TCP]
Antes de trocar qualquer dado, as duas pontas trocam três mensagens para
\textbf{sincronizar números de sequência} --- o mecanismo que depois
permite ao TCP detectar perda e reordenar:
\begin{enumerate}
\item \textbf{SYN}: o cliente propõe um número de sequência inicial.
\item \textbf{SYN-ACK}: o servidor confirma o número do cliente e propõe o
próprio.
\item \textbf{ACK}: o cliente confirma o número do servidor.
\end{enumerate}
Depois desse acordo, as duas pontas sabem contar exatamente quais bytes já
chegaram e quais faltam --- é isso que viabiliza a garantia de entrega e
ordem. O handshake \textbf{não autentica} identidade, \textbf{não negocia}
criptografia, e \textbf{não controla fluxo} de vazão --- essas são outras
funções (a autenticação/criptografia fica a cargo de camadas superiores
como TLS; o controle de fluxo é outro mecanismo do TCP, separado do
handshake). E só existe no TCP: o UDP não tem handshake porque não
estabelece conexão.
\end{conceito}

\section{Endereçamento e roteamento IP}

\subsection{IPv4, IPv6 e faixas privadas}

\begin{conceito}[Tamanho do endereço e por que existem faixas privadas]
\textbf{IPv4}: endereço de \textbf{32 bits}, escrito em 4 octetos decimais
(0--255). O espaço de endereços é finito (cerca de 4,3 bilhões), e não é
suficiente para dar um endereço público único a cada dispositivo do mundo
--- daí a RFC 1918 reservar faixas \textbf{privadas}, que podem ser
reutilizadas livremente dentro de redes internas diferentes (porque nunca
saem para a Internet pública diretamente): \texttt{10.0.0.0/8},
\texttt{172.16.0.0/12}, \texttt{192.168.0.0/16}. O NAT (adiante) é o que
permite esses endereços privados acessarem a Internet pública.

\textbf{IPv6}: endereço de \textbf{128 bits}, criado justamente para
resolver a escassez do IPv4 --- o espaço de endereços é astronomicamente
maior, o que elimina a necessidade estrutural de NAT (ainda que ele possa
existir por outros motivos). Prefixo de sub-rede típico: \texttt{/64}.
\end{conceito}

\begin{conceito}[Notação simplificada do IPv6]
Um endereço IPv6 completo tem 8 grupos de 4 dígitos hexadecimais. Duas
regras de abreviação, aplicadas \textbf{nesta ordem}:
\begin{enumerate}
\item Apague os \textbf{zeros à esquerda} de cada grupo individualmente
(\texttt{0DB8} $\to$ \texttt{DB8}; \texttt{0001} $\to$ \texttt{1}; um grupo
totalmente zero, \texttt{0000}, vira só \texttt{0}).
\item Substitua \textbf{uma única} sequência contígua de grupos
inteiramente nulos por \texttt{::} --- só pode ser usada \textbf{uma vez}
no endereço inteiro, porque usá-la duas vezes tornaria impossível saber
quantos grupos cada \texttt{::} está escondendo.
\end{enumerate}
\end{conceito}

\begin{exemplo}[Abreviando um endereço IPv6]
Endereço completo: \texttt{2001:0DB8:0000:0000:0000:0000:FE00:0001}.

Passo 1 (zeros à esquerda de cada grupo): \texttt{2001:DB8:0:0:0:0:FE00:1}.

Passo 2 (a maior sequência de grupos zerados vira \texttt{::} --- aqui, os
quatro grupos \texttt{0:0:0:0} no meio): \texttt{2001:DB8::FE00:1}.
\end{exemplo}

\subsection{Sub-redes: calculando hosts utilizáveis}

\begin{conceito}[Por que se subtrai 2 do total de endereços]
Um endereço IPv4 tem 32 bits, divididos entre a parte de \textbf{rede} e a
parte de \textbf{host}. O prefixo \texttt{/n} informa que os primeiros $n$
bits identificam a rede, sobrando $h = 32 - n$ bits para identificar hosts
dentro dela. Com $h$ bits, existem $2^h$ combinações possíveis --- mas duas
delas são \textbf{reservadas} e não podem ser atribuídas a um host:

\begin{itemize}
\item o endereço com todos os bits de host em \textbf{0} identifica a
\textbf{própria rede} (não um host);
\item o endereço com todos os bits de host em \textbf{1} é o
\textbf{broadcast} da sub-rede (não um host).
\end{itemize}

Por isso, \textbf{hosts utilizáveis $= 2^h - 2$}.
\end{conceito}

\begin{exemplo}[Aplicando a fórmula]
\begin{center}
\begin{tabular}{@{}p{1.6cm}p{3.6cm}p{2.2cm}p{3.4cm}@{}}
\toprule
\textbf{Prefixo} & \textbf{Máscara} & \textbf{Bits de host ($h$)} & \textbf{Hosts utilizáveis ($2^h{-}2$)} \\
\midrule
/24 & 255.255.255.0 & 8 & $2^8-2=254$ \\
/25 & 255.255.255.128 & 7 & $2^7-2=126$ \\
/26 & 255.255.255.192 & 6 & $2^6-2=62$ \\
/27 & 255.255.255.224 & 5 & $2^5-2=30$ \\
/28 & 255.255.255.240 & 4 & $2^4-2=14$ \\
\bottomrule
\end{tabular}
\end{center}
\end{exemplo}

\begin{regrapratica}[Do requisito de hosts para a máscara certa]
Quando o enunciado dá uma necessidade (“a sub-rede precisa comportar até 30
hosts”) e pede a máscara, procure a \textbf{menor} sub-rede (o maior
prefixo \texttt{/n}) que ainda \textbf{comporta} o número pedido --- usar
uma sub-rede maior que o necessário desperdiça endereços. Trinta hosts
cabem exatamente em \texttt{/27} (30 utilizáveis); \texttt{/28} (14) não
cabe; \texttt{/26} (62) cabe mas desperdiça 32 endereços à toa.
\end{regrapratica}

\subsection{Roteamento e distância administrativa}

\begin{conceito}[Rota estática $\times$ dinâmica, e como o roteador desempata]
\textbf{Rota estática}: configurada manualmente pelo administrador, fixa até
alguém mudar. \textbf{Rota dinâmica}: aprendida automaticamente por um
protocolo de roteamento que troca informação com outros roteadores (OSPF,
RIP, BGP) --- se adapta sozinha a mudanças na topologia.

Quando o roteador aprende \textbf{mais de uma rota} para o mesmo destino,
por fontes diferentes, ele precisa de um critério de desempate: a
\textbf{distância administrativa} (um número que representa a
“confiabilidade” atribuída a cada fonte de rota, menor = mais confiável).
Valores padrão (Cisco): rede diretamente conectada = 0, rota estática = 1,
OSPF = 110, RIP = 120. \textbf{A menor distância administrativa vence},
sempre --- é um critério puramente numérico, sem exceção por “protocolo
mais moderno” ou situação especial.
\end{conceito}

\begin{cuidado}[Não existe prioridade “inteligente” no roteamento]
Alternativas que descrevem comportamento adaptativo além da distância
administrativa --- “OSPF sempre tem prioridade sobre rota estática”,
“o roteador faz balanceamento de carga automático entre rotas
concorrentes”, “descarta pacotes quando há rotas conflitantes” --- estão
inventando comportamento. A regra real é mecânica e determinística: compara
a distância administrativa das rotas disponíveis para o mesmo destino e
usa a menor. Ponto.
\end{cuidado}

\subsection{ACL Cisco e \textit{wildcard mask}}

\begin{conceito}[Wildcard mask é o inverso lógico da máscara de sub-rede]
Uma \textit{wildcard mask} usa a lógica oposta da máscara normal: bit
\textbf{0} significa “esse bit \textbf{tem} que bater exatamente” e bit
\textbf{1} significa “esse bit é \textbf{ignorado}” (pode ser qualquer
valor). É obtida invertendo cada bit da máscara de sub-rede correspondente.
\end{conceito}

\begin{exemplo}[Construindo a wildcard mask para “só endereços ímpares”]
Um endereço IPv4 é ímpar quando o \textbf{bit menos significativo do último
octeto} vale 1. Para uma ACL casar \emph{apenas} endereços ímpares, esse
bit específico precisa estar \textbf{fixo} (wildcard = 0 nesse bit), e
todos os demais bits do último octeto podem ser qualquer coisa (wildcard =
1). Isso dá, no último octeto, \texttt{11111110} em binário $=$
\texttt{254} em decimal. Wildcard resultante: \texttt{0.0.0.254}. O
raciocínio é sempre \textbf{bit a bit}: monte o octeto em binário decidindo
“fixo (0)” ou “livre (1)” posição por posição, e só depois converta para
decimal --- decorar valores prontos leva a erro quando o padrão muda.
\end{exemplo}

\subsection{MPLS, VLAN e NAT}

\begin{conceito}[Três mecanismos que resolvem problemas diferentes na rede]
\begin{itemize}
\item \textbf{MPLS} (\textit{Multiprotocol Label Switching}): em vez de
cada roteador olhar o endereço IP de destino e recalcular a rota a cada
salto, o MPLS anexa um \textbf{rótulo} (\textit{label}) simples ao pacote
na entrada da rede, e os roteadores internos só olham o rótulo para
encaminhar --- muito mais rápido que analisar IP a cada salto. Por operar
“entre” a camada de enlace e a de rede (usa rótulos como a camada 2, mas
serve para encaminhar como a camada 3), é apelidado de \textbf{“camada
2.5”}. Permite criar caminhos dedicados (LSP) para engenharia de tráfego e
qualidade de serviço (QoS).
\item \textbf{VLAN} (\textit{Virtual LAN}): segmenta \textbf{logicamente} a
camada 2 --- cria domínios de \textbf{broadcast} separados dentro do
\textbf{mesmo} switch físico, sem precisar de switches físicos distintos
para cada grupo de máquinas. Um pacote de broadcast de uma VLAN não vaza
para outra, mesmo compartilhando o mesmo hardware.
\item \textbf{NAT} (\textit{Network Address Translation}): traduz
endereços \textbf{privados} (dentro da rede interna) para um endereço
\textbf{público} (ao sair para a Internet) e vice-versa na volta ---
resolve dois problemas ao mesmo tempo: economiza endereços IPv4 públicos
(muitos hosts internos compartilham um só IP público) e oculta a topologia
interna da rede de quem está de fora.
\end{itemize}
\end{conceito}

\begin{cuidado}[MPLS não é “roteamento IP mais rápido”; VLAN não é segmentação física]
MPLS não substitui o roteamento IP --- ele acelera o encaminhamento
\textbf{depois} que a rota já foi decidida, usando rótulo em vez de
recalcular IP a cada salto; não é a mesma coisa que “um protocolo de
roteamento melhor”. VLAN segmenta a camada \textbf{2} (domínio de
broadcast) por configuração lógica --- não é uma separação de cabos ou
hardware, é a mesma infraestrutura física dividida por software.
\end{cuidado}

\section{Tipos de rede corporativa}

\begin{conceito}[Internet, intranet, extranet e portal]
Os quatro termos descrevem \textbf{quem tem acesso} a uma rede, não sua
tecnologia (todos rodam sobre a mesma pilha TCP/IP):
\begin{itemize}
\item \textbf{Internet}: rede \textbf{pública} e global, acessível a
qualquer um.
\item \textbf{Intranet}: rede \textbf{interna} de uma organização, acesso
restrito aos seus membros.
\item \textbf{Extranet}: extensão controlada da intranet para
\textbf{parceiros ou clientes externos específicos} --- não é pública como
a Internet, mas também não é fechada só aos funcionários como a intranet;
é um meio-termo de acesso seletivo.
\item \textbf{Portal}: um \textbf{ponto único de acesso} organizado a
serviços/informações (pode ser a interface tanto de uma intranet quanto de
uma extranet).
\end{itemize}
\end{conceito}

\section{Segurança de rede}

\subsection{Firewall: \textit{stateless} $\times$ \textit{stateful}}

\begin{conceito}[A diferença é lembrar (ou não) do que já aconteceu]
\begin{center}
\begin{tabular}{@{}p{2.4cm}p{4.6cm}p{4.6cm}@{}}
\toprule
& \textbf{Stateless} (filtro de pacotes) & \textbf{Stateful} (com estado) \\
\midrule
O que examina & cada pacote \textbf{isoladamente} --- IP origem/destino,
porta, protocolo & o pacote \textbf{no contexto da conexão/sessão} a que
pertence \\
Como decide & aplica a regra estática ao pacote, sem memória do que passou
antes & mantém uma \textbf{tabela de estado} das conexões ativas e usa esse
histórico na decisão \\
Consequência prática & para permitir a resposta de uma conexão, é preciso
escrever uma regra explícita de volta & a resposta de uma conexão iniciada
de dentro é aceita \textbf{automaticamente}, porque o firewall já sabe que
ela pertence a uma sessão legítima \\
\bottomrule
\end{tabular}
\end{center}
Ambos operam nas camadas 3 e 4 (IP e porta). Quem inspeciona o
\textbf{conteúdo} da camada de aplicação (camada 7) é outra categoria de
equipamento --- o \textbf{firewall de próxima geração} (NGFW) ou o
\textbf{WAF} (\textit{Web Application Firewall}).
\end{conceito}

\begin{cuidado}[Não misture “com estado” com “opera na aplicação”]
Dizer que o firewall \textit{stateful} “atua exclusivamente na camada de
aplicação” é um absoluto errado --- ele continua trabalhando nas camadas
3/4, só que \textbf{lembrando} do contexto da conexão. Subir até a camada 7
é uma capacidade \textbf{adicional} de outra categoria de equipamento
(NGFW/WAF), não uma consequência automática de ser \textit{stateful}.
\end{cuidado}

\subsection{SSH substituindo o Telnet}

\begin{conceito}[A mesma função, com e sem cifra]
Telnet e SSH resolvem o mesmo problema --- \textbf{administração remota de
outro equipamento pela rede} --- mas com uma diferença que os tornou
incompatíveis em qualquer ambiente que leve segurança a sério:

\begin{itemize}
\item \textbf{Telnet} (porta 23): transmite \textbf{tudo em texto claro},
inclusive o usuário e a senha do administrador. Qualquer um capaz de
capturar o tráfego no caminho (um \textit{sniffer} na mesma rede, por
exemplo) lê a sessão inteira, senha incluída.
\item \textbf{SSH} (porta 22): faz o mesmo trabalho, mas \textbf{cifra toda
a comunicação}, incluindo a etapa de autenticação --- ninguém capturando o
tráfego consegue ler usuário, senha ou comandos.
\end{itemize}

A razão de o SSH ter substituído o Telnet é estritamente \textbf{segurança
(criptografia)}, não velocidade ou qualquer outra vantagem operacional.
\end{conceito}

\subsection{TLS: por que a criptografia é híbrida}

\begin{conceito}[Assimétrica só para combinar a chave; simétrica para o volume]
O TLS (o que torna HTTP em HTTPS) resolve dois problemas que pedem
ferramentas diferentes, e por isso combina os dois tipos de criptografia:

\begin{itemize}
\item \textbf{Criptografia assimétrica --- só no \textit{handshake}}: o
certificado digital do servidor comprova sua identidade (é ele mesmo quem
diz ser), e o par de chaves pública/privada permite às duas pontas
\textbf{combinarem uma chave simétrica de sessão} \textbf{sem que ela
trafegue em claro} pela rede.
\item \textbf{Criptografia simétrica --- para todo o resto}: uma vez
combinada a chave de sessão, \textbf{todo} o volume de dados trocado
durante a conexão é cifrado com ela, porque criptografia simétrica é
\textbf{ordens de magnitude mais rápida} que assimétrica.
\end{itemize}

O motivo do híbrido: a assimétrica sozinha seria lenta demais para cifrar
todo o tráfego de uma sessão inteira; a simétrica sozinha não resolveria o
problema de \emph{como} as duas pontas combinam uma chave secreta por um
canal que, até aquele momento, ainda não é seguro.
\end{conceito}

\begin{cuidado}[O que o handshake NÃO faz]
A parte assimétrica do TLS \textbf{não} cifra o tráfego inteiro da sessão
(só participa da negociação da chave) e \textbf{não} autentica o usuário
final por login/senha --- ela autentica o \textbf{servidor}, via
certificado digital. Autenticação do usuário (se houver) é uma camada
posterior, dentro da aplicação, não parte do handshake TLS.
\end{cuidado}

\begin{cuidado}[SSL não é “mais moderno” que TLS]
TLS é o \textbf{sucessor} do SSL, criado justamente para corrigir
vulnerabilidades encontradas no protocolo antigo. Qualquer alternativa que
descreva o SSL como mais seguro, mais recente ou mais recomendado que o TLS
inverteu a linha do tempo.
\end{cuidado}

\section{Roteiro de revisão do capítulo}

\begin{regrapratica}[Sequência de raciocínio para qualquer questão de redes]
\begin{enumerate}
\item O enunciado descreve uma \textbf{função}? Identifique a camada OSI
pelo comportamento (Seção anterior), não tente decorar “camada X = tal
coisa” sem entender o porquê.
\item Envolve \textbf{dois hosts trocando dados}? Pergunte se é TCP
(precisa de confiabilidade/ordem) ou UDP (velocidade importa mais).
\item Envolve \textbf{endereçamento}? Verifique se é cálculo de sub-rede
($2^h - 2$), faixa privada (RFC 1918) ou notação IPv6.
\item Envolve \textbf{decisão de caminho}? Verifique se é rota
estática/dinâmica e aplique a distância administrativa (menor vence).
\item Envolve \textbf{segurança}? Separe o que é do protocolo de
transporte/rede (TCP handshake, firewall stateful) do que é de
criptografia de aplicação (TLS, SSH) --- são mecanismos independentes que a
prova gosta de combinar no mesmo cenário.
\end{enumerate}
\end{regrapratica}
